% =====================================================================
%  05_datos.tex — Padrón, esquema SQLite, vistas y consultas
% =====================================================================
\chapter{Datos y métricas}
\label{cap:datos}

\section{El padrón}

La única fuente de las unidades válidas es \path{padron.csv}, con tres
columnas:

\begin{lstlisting}[caption={Formato de \texttt{padron.csv}}, label={lst:padron}]
no_economico,empresa,tipo
5038,CISA,electrico
5014,CISA,diesel
\end{lstlisting}

Reglas que aplica \path{config.py} al leerlo:

\begin{itemize}
  \item \texttt{no\_economico} debe ser numérico y no repetirse.
  \item \texttt{empresa} no puede estar vacía.
  \item \texttt{tipo} debe ser \texttt{electrica} o \texttt{diesel}.
  \item Si una fila está mal, el sistema no arranca y dice el número de línea.
\end{itemize}

Al iniciar, \path{metricas/storage.py} sincroniza la tabla \texttt{unidad} con
el archivo. Para dar de alta o baja una unidad: edita el CSV y reinicia. Los
números detectados que no están en el padrón también se guardan en
\texttt{unidad}, con \texttt{en\_padron = 0} y empresa
\texttt{Desconocida}.

\section{Esquema de la base}

La base es SQLite (\path{datos/detecciones.db} dentro de Docker) y usa el modo
WAL. El esquema vive en \path{metricas/schema.sql} y es el \textbf{único} del
sistema: se aplica en cada arranque con \texttt{CREATE ... IF NOT EXISTS}.

\begin{longtable}{@{}L{3.6cm}L{10.6cm}@{}}
  \toprule
  \textbf{Tabla} & \textbf{Contenido} \\
  \midrule
  \endhead
  \texttt{tipo\_unidad} & Catálogo: \texttt{electrica}, \texttt{diesel},
    \texttt{desconocido}. \\
  \texttt{puerta} & Catálogo de puertas del patio (\texttt{uno},
    \texttt{dos}). \\
  \texttt{unidad} & Una fila por número económico, con empresa, tipo y
    \texttt{en\_padron}. Se sincroniza desde el CSV. \\
  \texttt{evento\_paso} & Una fila por cruce: unidad, dirección, puerta, zona,
    confianza, foto y hora del paso. \\
  \texttt{estancia} & Una fila por visita: se abre al entrar y se cierra al
    salir. \\
  \bottomrule
  \caption{Tablas de la base de datos.}
  \label{tab:tablas}
\end{longtable}

\subsection{\texttt{evento\_paso}}

\begin{center}
\begin{tabular}{@{}llp{7.6cm}@{}}
  \toprule
  \textbf{Columna} & \textbf{Tipo} & \textbf{Significado} \\
  \midrule
  \texttt{id\_evento}     & entero  & Llave primaria. \\
  \texttt{id\_unidad}     & entero  & Unidad que cruzó. \\
  \texttt{direccion}      & texto   & \texttt{entrada} o \texttt{salida}
                                        (validado con \texttt{CHECK}). \\
  \texttt{id\_puerta}     & entero  & Puerta donde ocurrió. \\
  \texttt{zona}           & texto   & \texttt{delantero}, \texttt{lateral} o
                                        \texttt{trasera}. \\
  \texttt{confianza}      & real    & Entre 0 y 1. \\
  \texttt{ruta\_captura}  & texto   & Ruta de la foto del aviso. \\
  \texttt{hora\_paso}     & texto   & Hora local del patio,
                                        \texttt{AAAA-MM-DDTHH:MM:SS}. \\
  \texttt{hora\_registro} & texto   & Hora en que se escribió la fila. \\
  \bottomrule
\end{tabular}
\end{center}

Índices: por \texttt{hora\_paso} descendente y por
(\texttt{id\_unidad}, \texttt{hora\_paso}).

\subsection{\texttt{estancia}}

Una estancia resume un par entrada/salida. La duración y el estatus
\textbf{no se escriben}: son columnas generadas por SQLite.

\begin{lstlisting}[language=SQL, caption={Columnas generadas de \texttt{estancia}},
                   label={lst:estancia}]
minutos REAL GENERATED ALWAYS AS (
  CASE WHEN hora_salida IS NOT NULL
  THEN ROUND((julianday(hora_salida) - julianday(hora_entrada)) * 1440.0, 1)
  END
) STORED,
estatus TEXT GENERATED ALWAYS AS (
  CASE WHEN hora_salida IS NULL THEN 'dentro' ELSE 'fuera' END
) VIRTUAL
\end{lstlisting}

Mientras \texttt{hora\_salida} es \texttt{NULL}, la unidad sigue adentro y su
estatus es \texttt{dentro}.

\section{Vistas}

El código consulta vistas, no arma \texttt{JOIN} a mano:

\begin{longtable}{@{}L{4.6cm}L{9.6cm}@{}}
  \toprule
  \textbf{Vista} & \textbf{Devuelve} \\
  \midrule
  \endhead
  \texttt{v\_estancia} & Estancias con número económico y empresa. \\
  \texttt{v\_evento} & Eventos con empresa, tipo de unidad y bandera de
    padrón. \\
  \texttt{v\_dentro\_ahora} & Unidades que están dentro en este momento, con
    minutos transcurridos. \\
  \texttt{v\_historial\_estancias} & Historial ordenado por hora de entrada. \\
  \texttt{v\_promedio\_estancia} & Visitas, promedio, mínimo y máximo por
    empresa. \\
  \texttt{v\_no\_registrados} & Números fuera del padrón, con veces y última
    aparición. \\
  \bottomrule
  \caption{Vistas de la base de datos.}
  \label{tab:vistas}
\end{longtable}

\section{Reglas de registro}

\begin{enumerate}
  \item \textbf{Un evento por cruce.} El registro se hace al terminar de pasar
        el camión, no al tocar la línea.
  \item \textbf{Debounce.} El mismo número no se repite antes de
        \texttt{DEBOUNCE\_SEG} (60 s).
  \item \textbf{Maniobra.} Si la unidad regresa antes de \texttt{MANIOBRA\_SEG}
        (20 s), no se registra.
  \item \textbf{Día operativo.} El día empieza a las \textbf{3:00 a.m.}
        (\texttt{HORA\_CORTE\_DIA}): los movimientos de la madrugada cuentan
        para el día anterior.
  \item \textbf{Fuera del padrón.} Con \texttt{VOTOS\_FUERA\_PADRON} lecturas
        iguales se registra el número tal cual, con \texttt{en\_padron = 0}.
\end{enumerate}

\section{Consultas útiles}

\begin{lstlisting}[language=SQL, caption={Unidades dentro del patio ahora},
                   label={lst:q-dentro}]
SELECT no_economico, empresa, hora_entrada, minutos_dentro
FROM v_dentro_ahora;
\end{lstlisting}

\begin{lstlisting}[language=SQL, caption={Movimientos de hoy},
                   label={lst:q-hoy}]
SELECT no_economico, empresa, direccion, zona, hora_paso
FROM v_evento
WHERE date(hora_paso) = date('now', 'localtime')
ORDER BY hora_paso DESC;
\end{lstlisting}

\begin{lstlisting}[language=SQL, caption={Números detectados fuera del padrón},
                   label={lst:q-fuera}]
SELECT no_economico, veces, ultima_vez
FROM v_no_registrados
ORDER BY veces DESC;
\end{lstlisting}

\begin{lstlisting}[language=SQL, caption={Promedio de estancia por empresa},
                   label={lst:q-promedio}]
SELECT empresa, total_visitas, promedio_minutos, minimo_minutos, maximo_minutos
FROM v_promedio_estancia;
\end{lstlisting}

\section{Respaldo y migración}

\begin{itemize}
  \item \textbf{Respaldo}: copia \path{datos/detecciones.db} (y, si existe,
        \texttt{detecciones.db-wal}). Con el sistema detenido el respaldo es
        consistente.
  \item \textbf{Migración}: si \path{storage.py} detecta el formato anterior,
        respalda la base y la convierte al esquema actual antes de seguir.
  \item \textbf{No edites la base a mano} mientras el sistema corre: usa las
        vistas o detén los servicios primero.
\end{itemize}
